\section{Measurement accuracy and causal validity}
\label{reading:assistance-conclusions}

\subsection{Better measurement can still leave confounding}

Extracting financial distress from customer notes illustrates the distinction.
A model may summarize the notes accurately while the notes omit circumstances
that influenced both a servicing action and repayment. The measurement example
in Section~\ref{exec13:census-measurement} makes the consequence explicit.
With true distress $U$, measured distress $W=U+\varepsilon_W$, action
$A=U+\varepsilon_A$ and outcome $Y=\tau A+U+\varepsilon_Y$, mutually
independent centred Gaussian disturbances yield
\[
 \beta_{A\mid W}=\tau+\frac{s}{1+2s},\qquad
 s=\operatorname{Var}(\varepsilon_W),
\]
when the other variances are one. The second term is residual confounding.
It persists in an arbitrarily large sample. A narrower statistical interval
around the adjusted coefficient can therefore make an incorrect economic
interpretation appear more certain. The primary review of text as confounder
measurement discusses the need to distinguish confounders, mediators and
colliders; the displayed calculation is the manuscript's constructed example.
\citep{exec13_census_text_review_arxiv}

\begin{figure}[htbp]\centering
\begin{tikzpicture}[x=2.2cm,y=6cm,>=Stealth,every node/.style={font=\small}]
 \draw[->] (0,-.25)--(0,.3) node[above]{Effect units};
 \draw[->] (0,-.25)--(4.35,-.25) node[right]{Measurement-noise variance};
 \draw[dashed,gray] (0,0)--(4.1,0);
 \draw[thick,ink,dashed] (0,-.2)--(4,-.2);
 \draw[very thick,teal,domain=0:4,samples=60] plot (\x, {-.2+\x/(1+2*\x)});
 \foreach \x in {1,2,3,4}{\draw (\x,-.26)--(\x,-.24);\node[below] at (\x,-.27){\x};}
 \node[left] at (0,-.2){$-0.2$};\node[left] at (0,0){$0$};\node[left] at (0,.2){$0.2$};
 \node[above,align=center] at (2.7,.23){Adjusted regression\ coefficient};
 \node[above] at (2.7,-.2){Specified causal effect};
\end{tikzpicture}
\caption{In the constructed Gaussian example, measurement error can reverse
the sign of an adjusted regression coefficient. The causal effect is fixed at
$\tau=-1/5$; the solid curve is $-1/5+s/(1+2s)$. This is an exact population
calculation, not observed performance of a language model.}
\label{reading:measurement-error}
\end{figure}

The appropriate comparison is consequently broader than extraction accuracy.
Simple dictionaries, supervised extraction, human coding and language models
should use the same permitted document dates and the same causal target.
The evaluation should ask whether their measurement errors change the estimated
effect or the chosen action. It should also test cases in which the notes were
written after treatment or only after a complaint. Those cases alter the
causal role of the text, not merely its predictive quality.
\citep{exec13_census_text_review_arxiv}

This example concerns the validity of the measurement used for adjustment.
The task comparisons in Chapter~\ref{rewrite:assistance-tasks} separately
address formulation, discovery, formal reasoning, execution and reporting.
They distinguish supplied-structure benchmarks from elicitation of a new
bank mechanism. \citep{v12_la_2312_04350,v12_la_2606_23608}
